\section{Results}
\label{sec:results}

The estimation pipeline of \Model operates as follows: given an input AFM image, the trained g-CNN directly predicts the CV values (inter-domain distances) from the image; these predicted CV values are then used as restraints to navigate the generative process of \AFiii, which produces a 3D conformation consistent with the input AFM image.
Because the g-CNN is invariant to in-plane rotations and translations, the same conformation yields the same CV estimate regardless of how the molecule is oriented or positioned in the image, eliminating the need for explicit image alignment.
Note that, because the CVs and the training conformational ensemble are specific to each target protein, the g-CNN must be trained separately for each protein system of interest.

\subsection{Verification using pseudo-AFM images: adenylate kinase}
\label{subsec:ak_results}

Adenylate kinase (\AK) is a well-studied monomeric enzyme known for its large-scale conformational transitions. It is composed of three relatively rigid domains: the central CORE domain (residues 1–29, 68–117, and 161–214), the AMP-binding domain (AMPbd; residues 30–67), and the lid-like ATP-binding domain (LID; residues 118–167). Experimental and computational studies have suggested that upon ligand binding, the enzyme undergoes a transition from an inactive open conformation to an active closed conformation (see \Cref{fig:overview}~(c)). Here, we evaluated how accurately \Model reconstructs 3D conformations of \AK from artificially created pseudo ``experimental'' AFM images. 

\begin{figure}[htbp]
    \centering
    \includegraphics[width=0.75\linewidth]{images/figure-2.png}
    \caption{
        \textbf{Training data and estimation accuracy of \Model}.
        \textbf{(a)} Projection of training data into the inter-domain distance space. The two cross marks represent the open (4AKE) and closed (1AKE) crystal structures from PDB, respectively.
        \textbf{(b)} Projection of a 450 ns MD simulation trajectory into the inter-domain distance space.
        \textbf{(c)} Distribution of the root squared error in the inter-domain distance space between the g-CNN's estimations and the ground-truth. The inter-domain distance was estimated by the trained g-CNN from pseudo-AFM images, which were generated from randomly selected molecular dynamics simulation conformations. The average of the root squared errors is indicated by the broken line (0.254 nm). 
        \textbf{(d)} Scatter plot of the root squared error in the inter-domain distance space and the heavy-atom RMSDs from the ground-truth conformation. 
        The heavy-atom RMSD was computed between the conformation generated with \AFiii navigating using the predicted CV values, and the ground-truth conformation.
        \textbf{(e)} Density of the heavy-atom RMSDs. The average of the RMSD is indicated by the broken line (0.281 nm).
        }
    \label{fig:ak_cnn_training}
\end{figure}

\paragraph{Training data generation and g-CNN accuracy}
We generated a conformational ensemble by navigating \AFiii to target collective variable (CV) values---defined as the inter-domain distances for all domain pairs (LID--CORE, CORE--AMPbd, and AMPbd--LID)---
sampled on a regular grid (see SI section S1 for details of the navigation schedule).
The generated ensembles are shown with their MolProbity scores in \Cref{fig:ak_selection}.
The resulting ensemble adequately covers the broad conformational space of \AK, encompassing both the closed (PDB ID: 1AKE) and open (PDB ID: 4AKE) crystal structures (\Cref{fig:ak_cnn_training}~(a,b)).
We then trained the g-CNN on this ensemble and evaluated its prediction accuracy using pseudo-AFM images generated from randomly selected MD conformations (see \Cref{tab:afm_settings} for the image settings).
The mean root squared error in the CV space is 0.254~nm (\Cref{fig:ak_cnn_training}~(c)), apart from a few outliers exceeding 0.500~nm---sufficiently small to resolve the intermediate conformational states between the open and closed forms of \AK.

\paragraph{Accuracy of conformations estimated by \Model}
We first applied \Model to pseudo-AFM images of the closed and open crystal structures of \AK (\Cref{fig:overview}~(b,c)).
Note that these pseudo-AFM images were generated under idealized imaging conditions (small tip radius, fine resolution, no noise; see Methods), so the accuracy reported here represents an upper bound on what can be achieved; the impact of more realistic conditions is assessed in the robustness analysis below.
The RMSDs between the reconstructed and ground-truth conformations were 0.216~nm (closed) and 0.176~nm (open), both well below the 0.715~nm RMSD separating the two crystal structures, confirming that \Model clearly distinguishes between the two states.
Accordingly, the pseudo-AFM images generated from the estimated conformations are visually similar from the input images, with correlation coefficients (c.c.; \Cref{eq:cc}) of 0.997 (closed) and 0.998 (open).

We next assessed reconstruction accuracy over a broader range of conformations using pseudo-AFM images generated from randomly selected MD snapshots.
Note that these MD conformations span a continuous distribution across the open--closed landscape and are not classified into discrete states; the RMSD and CV errors reported below therefore reflect accuracy over this continuum.
Most RMSDs are distributed around the mean of 0.281~nm, with a small number of outliers (\Cref{fig:ak_cnn_training}~(d,e)).
Inspection of the outlier cases revealed two contributing factors (\Cref{fig:ak_cnn_training}~(d)).
First, depending on how the molecule is oriented on the stage (i.e., which face of the protein contacts the surface), the resulting AFM image may look similar regardless of whether the conformation is open or closed.
Because the g-CNN is trained on pseudo-AFM images rendered from uniformly sampled 3D orientations, it must learn to predict CVs from all possible views; however, certain out-of-plane orientations produce nearly degenerate images for structurally distinct conformations, making it inherently difficult for the g-CNN to estimate the correct inter-domain distances from such views.
Second, even when the predicted CV values are accurate, the reconstructed RMSD can still be elevated because the CV--structure mapping is not unique: 
distinct domain arrangements---such as hinge-bending versus shear-like motions---can yield similar inter-domain distances while producing structurally different conformations.
Indeed, \Cref{fig:ak_cnn_training}~(d) shows that several data points with small CV errors nonetheless exhibit elevated RMSDs, confirming that the CV--structure degeneracy---rather than CV prediction inaccuracy---is the dominant source of reconstruction error in these cases.

\begin{figure}[htbp]
    \centering
    \includegraphics[width=0.75\linewidth]{images/figure-3.png}
    \caption{
        \textbf{Robustness against noise and imaging conditions}.
        \textbf{(a)} Examples of pseudo-AFM images with Gaussian noise at various standard deviations (0.1--0.4~nm), generated from randomly selected MD conformations.
        \textbf{(b)} Mean RMSD between the estimated and ground-truth structures as a function of noise levels in both the training and ``experimental'' pseudo-AFM images. 
        Each value was computed from 20 independent estimations.
        \textbf{(c)} Mean squared error (MSE) of the inter-domain distances under the same noise conditions as in (b).
        \textbf{(d)} Pseudo-AFM images generated with different combinations of the image resolution $R_{\rm res}$ and the probe tip radius $R_{\rm probe}$.
        \textbf{(e)} Mean RMSD between the estimated and ground-truth structures as a function of the tip radius and the image resolution.
        \textbf{(f)} MSE of the inter-domain distances as a function of the tip radius and the image resolution.
    }
    \label{fig:noise_robustness}
\end{figure}

\paragraph{Robustness against noise and imaging conditions}
To evaluate the applicability of \Model to real experimental data, 
we investigated its robustness against two key factors that degrade image quality: 
measurement noise and imaging conditions (image resolution and probe tip radius).

We first examined the effect of noise by systematically varying the Gaussian noise levels added to both the training and ``experimental'' pseudo-AFM images (\Cref{fig:noise_robustness}~(a)).
The resulting heat maps of the mean RMSD (\Cref{fig:noise_robustness}~(b)) and the inter-domain distance MSE (\Cref{fig:noise_robustness}~(c)) reveal two distinct trends.
Adding noise to the ``experimental'' images consistently led to substantial reductions in prediction accuracy, whereas adding noise to the training data alone had a comparatively minor effect; 
indeed, the best performance was achieved with noise-free training data.
These results indicate that \Model is sensitive to noise in the input images but does not require noisy training data to handle noisy inputs.

We next assessed the influence of imaging conditions by varying the image resolution $R_{\rm res}$ and the probe tip radius $R_{\rm probe}$ (\Cref{fig:noise_robustness}~(d)).
The heat maps of the mean RMSD (\Cref{fig:noise_robustness}~(e)) and the inter-domain distance MSE (\Cref{fig:noise_robustness}~(f)) show that image resolution is the dominant factor governing prediction accuracy.
Because \AK is a relatively small molecule, coarse resolution causes it to appear as a nearly featureless round spot, making it difficult for the g-CNN to distinguish different conformations; 
accordingly, degrading the resolution leads to a dramatic increase in estimation error.
In contrast, increasing the probe tip radius has a comparatively modest effect on accuracy, 
indicating that \Model is relatively robust to tip broadening effects.

\paragraph{Correlation coefficient as an evaluation metric}
\begin{figure}[htbp]
    \centering
    \includegraphics[width=0.75\linewidth]{images/figure-4.png}
    \caption{
        \textbf{Variability of the correlation coefficient and its dependence on fitting parameters}.
        \textbf{(a)} Distribution of correlation coefficients (c.c.) between the ``experimental'' pseudo-AFM images and pseudo-AFM images generated from the estimated conformations after rigid-body fitting. The dashed line indicates the mean value (0.998).
        \textbf{(b)} Scatter plot of the c.c.\ versus the RMSD from the ground-truth conformation. The dashed line shows the linear fit.
        \textbf{(c--e)} Maximum c.c.\ achieved for different combinations of the translation grid step and the probe radius step at tip angles of $10^{\circ}$ (c), $20^{\circ}$ (d), and $30^{\circ}$ (e).
        \textbf{(f--h)} Maximum c.c.\ as a function of the number of rotation search steps at tip angles of $10^{\circ}$ (f), $20^{\circ}$ (g), and $30^{\circ}$ (h). Different curves correspond to different combinations of the probe radius step and translation grid step.
    }
    \label{fig:cc}
\end{figure}

When applying \Model to real experimental data, ground-truth structures are unavailable and RMSD cannot be computed. 
We therefore examined whether the correlation coefficient (c.c.) between pseudo-AFM images and the reference images can serve as a reliable proxy for structural accuracy.

The c.c.\ values are concentrated in a narrow range of 0.995--0.999, with a mean of 0.998 (\Cref{fig:cc}~(a)). 
Despite this narrow distribution, a clear negative correlation exists between the c.c.\ and the RMSD (\Cref{fig:cc}~(b)), confirming that the c.c.\ is indeed a meaningful indicator of structural agreement. 
Notably, however, even a modest difference of $\sim$0.01 in c.c.\ can correspond to a wide range of RMSD values (\Cref{fig:cc}~(b)), implying that small improvements in c.c.\ often reflect substantial differences in structural accuracy. 
This observation underscores the importance of precise c.c.\ evaluation: the metric is reliable, but the structurally relevant signal resides in very small numerical differences.

We next investigated how sensitively the c.c.\ depends on the parameters of the rigid-body fitting procedure. 
Using a single experimental AFM image and the crystal structure of \flhac (discussed in the next section) as the target, we performed repeated fitting trials in which the structure was randomly rotated and then optimized over translation and probe tip radius at various grid step sizes (\Cref{fig:cc}~(c--e)). 
The c.c.\ was highly sensitive to the translation grid step: a coarse translation search frequently failed to locate the optimal c.c.\ value. 
In contrast, the c.c.\ was relatively insensitive to the probe radius grid step. 
For the rotational search, the c.c.\ converged in most cases by approximately 5,000 rotation steps (\Cref{fig:cc}~(f--h)). 
Based on these results, we adopted a translation step of 0.7~nm, a probe radius step of 1~nm (2~nm -- 6~nm for \flhac), a tip angle of $20^{\circ}$, and 10,000 rotation steps for all subsequent analyses, which we consider sufficient for reliable c.c.\ evaluation.
With these settings, the rigid-body fitting takes approximately 3 minutes per structure on a single GPU.

\subsection{Application to real experimental AFM images: a flagellar protein \flhac}

FlhA is a key component of the flagellar protein export system, and its C-terminal domain, \flhac, helps transport proteins by assisting their binding to the exporter. 
Structurally, it comprises four domains (\acdi, \acdii, \acdiii, and \acdiv) and a linker to the transmembrane domain (FlhA\textsubscript{TM}) (\Cref{fig:flhac_results}~(a)). 
Previous studies using X-ray crystallography and mutational analysis indicate that hinge motions among these domains influence transport activity and motility \cite{saijo-hamano2010flha}. 
%Here, we applied \Model to real experimental AFM images of \flhac to evaluate its ability to reconstruct the 3D structure of \flhac from AFM images. 

Here, to evaluate the applicability of \Model to real AFM data, we used HS-AFM images of \flhac acquired by Inoue et al.\ \cite{inoue_structural_2019}.
The images were recorded with a spatial resolution of approximately 1.0~nm/pixel (59 $\times$ 59~nm$^2$ field of view, 60 $\times$ 60 pixels) and a frame interval of 66~ms, with the sample immobilized on mica in a buffer containing 20~mM NaCl and 10~mM Tris-HCl (pH~8.0).
We estimated the conformations of \flhac using \Model from 159 HS-AFM images of the \flhac monomer.
It should be noted that, whereas conventional flexible fitting methods for capturing structural changes typically require hours to days of computation per AFM image even with a coarse-grained model \cite{Niina2020},
\Model allows estimation of conformations from each AFM image in approximately one minute on a single GPU once the g-CNN has been trained (see Conclusion for a discussion of the training cost).
This enabled, for the first time, the processing of as many as 159 AFM images in this study.

\paragraph{Training of the g-CNN and ResNet-18}
Conformations for training the g-CNN were generated by navigating \AFiii with CV values on a grid as described in \textit{Training}. 
As the CVs, we used the inter-domain distances of domain pairs, \acdi--\acdiv, \acdii--\acdiii, and \acdii--\acdiv. 
The generated ensembles are shown with their MolProbity scores in \Cref{fig:flhac_selection}. 
The g-CNN was trained on the training data to learn the relationship between AFM images and the CVs.

We also trained ResNet-18 on the combination of the g-CNN's training images and the experimental AFM images 
to obtain embedding representations that capture the molecular features in AFM images.
To verify that the learned embeddings encode structurally meaningful information, 
we projected the 256-dimensional embedding vectors onto two dimensions using UMAP \cite{mcinnes2018umap} (\Cref{fig:resnet_embeddings}~(a--c)). 
The resulting scatter plots, colored by each of the three inter-domain distances, show that different inter-domain distances correspond to distinct directions in the embedding space, 
indicating that the embeddings successfully capture conformational variations.
We further confirmed that the embeddings are approximately invariant to in-plane rotation of the molecule.
For each of 20 randomly selected conformations, we generated pseudo-AFM images at $15^{\circ}$ rotation increments and computed their pairwise distances in the embedding space (\Cref{fig:resnet_embeddings}~(d)).
The cosine similarity between embeddings of rotationally related images was consistently high (mean = 0.948; \Cref{fig:resnet_embeddings}~(e)), 
and the L2 distance was consistently low (mean = 0.179; \Cref{fig:resnet_embeddings}~(f)), 
both well separated from the distributions of randomly sampled pairs.
These results, together with visual inspection, support the use of the embedding distance as a measure of image similarity: 
images that are close in the embedding space can be regarded as depicting similar molecular shapes.

\begin{figure}[htbp]
    \centering
    \includegraphics[width=0.75\linewidth]{images/figure-5.png}
    \caption{
        \textbf{Embedding representations learned by ResNet-18}.
        \textbf{(a--c)} UMAP projections of the 256-dimensional ResNet-18 embeddings for the training pseudo-AFM images and the experimental AFM images (red). Each scatter plot is colored by a different inter-domain distance: \acdii--\acdiv (a), \acdii--\acdiii (b), and \acdi--\acdiv (c). Different inter-domain distances correspond to distinct gradient directions in the embedding space.
        \textbf{(d)} Representative pseudo-AFM images of a single conformation rotated in the image plane at $15^{\circ}$ increments. Twenty such sets were prepared for statistical evaluation.
        \textbf{(e)} Distribution of cosine similarity in the embedding space. The gray histogram shows the distribution for randomly sampled image pairs from the training data. The blue dashed line indicates the mean cosine similarity for rotationally related image pairs (0.948), which remains high even when the variance is taken into account.
        \textbf{(f)} Distribution of L2 distance in the embedding space. The gray histogram shows the distribution for randomly sampled pairs. The green dashed line indicates the mean L2 distance for rotationally related pairs (0.179), which remains low even when the variance is taken into account.
    }
    \label{fig:resnet_embeddings}
\end{figure}

\paragraph{Evaluating estimated conformations by \Model}
We applied \Model to all 159 experimental AFM images and quantified the agreement between each estimated (or crystal) structure and the corresponding experimental image using two complementary metrics: 
the correlation coefficient (c.c.) and the L2 distance in the ResNet-18 embedding space. 
For each metric, we searched for the rendering pose that maximizes the similarity to the experimental image.

Both metrics consistently showed that \Model estimations agree with the experimental images better than the crystal structure.
In the embedding-space evaluation, the \Model estimations achieved smaller L2 distances than the crystal structure for the majority of frames (\Cref{fig:flhac_results}~(i)), with most distances falling below 0.3---within the rotational invariance range established in \Cref{fig:resnet_embeddings}~(f)---confirming that the generated images share similar shape features with the experimental data.
Likewise, in the c.c.-based evaluation, \Model achieved higher c.c.\ values in the majority of cases (\Cref{fig:flhac_results}~(j)).
\Cref{fig:flhac_results}~(b--f) illustrates a case in which the crystal structure clearly fails to reproduce the experimental image, whereas the \Model estimation captures the observed features in both the L2-based (\Cref{fig:flhac_results}~(c,d)) and c.c.-based (\Cref{fig:flhac_results}~(e,f)) fittings.

It should be noted, however, that the numerical margins are modest in both metrics. This reflects an intrinsic degeneracy in AFM imaging: distinct 3D conformations can yield similar AFM images due to the convolution effects of the probe tip, finite resolution, and measurement noise.
Indeed, although the embedding space analysis revealed that experimental images are predominantly surrounded by open conformations (\Cref{fig:flhac_results}~(g)), the closed crystal structure was still able to produce pseudo-AFM images of comparable similarity (\Cref{fig:flhac_results}~(h)), because under realistic imaging conditions, even a closed conformation can mimic the appearance of an open one.
Moreover, the c.c.\ is dominated by the overall molecular envelope and is inherently insensitive to fine internal structural rearrangements, as demonstrated for \AK (\Cref{fig:cc}~(a)); consequently, even meaningful conformational differences manifest as only small numerical changes in c.c.\ values.
Nevertheless, the fact that \Model consistently outperforms the crystal structure across both metrics and the majority of frames supports the conclusion that it captures genuine conformational features present in the experimental AFM images.

\begin{figure}[htbp]
    \centering
    \includegraphics[width=0.75\linewidth]{images/figure-6.png}
    \caption{
        \textbf{Comparison of \Model estimations and the crystal structure for \flhac}.
        \textbf{(a)} Crystal structure of \flhac (PDB ID: 3A5I). The four domains are colored as follows: dark blue, linker; sky blue, \acdi; green, \acdii; yellow, \acdiii; and orange, \acdiv.
        \textbf{(b)} A representative experimental AFM image in which the crystal structure poorly reproduces the observed features.
        \textbf{(c,d)} Pseudo-AFM images generated from the crystal structure (c) and the \Model estimation (d) after rigid-body fitting based on the L2 distance in the embedding space. Predicted domain positions are overlaid.
        \textbf{(e,f)} Pseudo-AFM images generated from the crystal structure (e) and the \Model estimation (f) after rigid-body fitting based on the c.c.\ value. Predicted domain positions are overlaid.
        \textbf{(g,h)} UMAP projections of the ResNet-18 embeddings, showing the positions of the fitted pseudo-AFM images (blue in (g), green in (h)) relative to the training data (gray) and experimental images (red), for the \Model estimations (g) and the crystal structure (h).
        \textbf{(i)} Scatter plot comparing the L2 distances from the experimental image to the \Model estimation ($y$-axis) and to the crystal structure ($x$-axis). Points below the diagonal indicate cases where \Model achieves a closer match. Marginal histograms on the top and right axes show the distributions of L2 distances for the crystal structure and \Model, respectively.
        \textbf{(j)} Bland--Altman plot of the c.c.\ values. The $x$-axis shows the mean of $\mathrm{c.c._{AFMFold}}$ and $\mathrm{c.c._{Crystal}}$, and the $y$-axis shows their difference. Points are colored by the \acdii--\acdiv inter-domain distance. The dashed line indicates the mean difference. The marginal histogram on the right shows the distribution of c.c.\ differences (improvement of \Model over the crystal structure).
    }
    \label{fig:flhac_results}
\end{figure}


\begin{figure}[htbp]
    \centering
    \includegraphics[width=0.75\linewidth]{images/figure-7.png}
    \caption{
        \textbf{Conformational dynamics of \flhac estimated by \Model}.
        \textbf{(a)} Distribution of the \acdii--\acdiv inter-domain distance estimated by \Model (red shaded) compared with AI-based structure generation models: AlphaFlow \cite{jing2024alphaflow}, BioEmu \cite{bioemu2025}, and \AFii with MSA subsampling \cite{delAlamo2022sampling}.
        \textbf{(b)} Time series of the six inter-domain distances estimated by \Model for the consecutive 159 experimental AFM frames.
        \textbf{(c)} Auto-correlation functions of the inter-domain distances as a function of time lag. Solid lines show the results from \Model estimations, while dashed lines correspond to temporally shuffled surrogate data.
    }
    \label{fig:flhac_statistic}
\end{figure}

\paragraph{Conformational dynamics}
To characterize the conformational dynamics captured by \Model, we computed the distribution of the \acdii--\acdiv inter-domain distance---a key indicator distinguishing between open and closed conformations---for all 159 estimated structures (\Cref{fig:flhac_statistic}~(a)).
For comparison, we also generated conformational ensembles using AI-based structure generation models: AlphaFlow \cite{jing2024alphaflow}, BioEmu \cite{bioemu2025}, and \AFii with MSA subsampling \cite{delAlamo2022sampling}, which are expected to approximate long-time behaviors of \flhac that are not accessible by brute-force MD simulations.
Interestingly, \Model tends to estimate more open conformational states than these AI-based models.
There are three possible explanations for this tendency.
First, HS-AFM captures molecular dynamics with a frame interval of 66~ms, which may reflect longer-timescale motions than those sampled by AlphaFlow or BioEmu, thereby enabling the observation of more open structures.
Second, interactions between the molecule and the stage during AFM imaging may stabilize conformations with larger contact surface areas, thus favoring more open structures.
Third, estimation errors in \Model could lead to an overrepresentation of highly open conformations with large conformational entropy.

To assess the third possibility, we examined the time series of the estimated inter-domain distances (\Cref{fig:flhac_statistic}~(b)).
Since the 159 AFM images are consecutive frames, we can expect temporal correlations in the estimated structures if the estimation errors are sufficiently small.
Indeed, the time series exhibits clear oscillatory behavior, with the distances fluctuating between larger (open) and smaller (closed) values over time.
Crucially, \Model estimates each frame independently without any temporal information, yet the estimated inter-domain distances display pronounced time correlations, implying that the estimation errors are smaller than the scale of the observed fluctuations.

This temporal correlation was quantified by the auto-correlation function (\Cref{fig:flhac_statistic}~(c)).
The auto-correlations from \Model estimations (solid lines) show a pronounced positive correlation for lag times up to 400~ms and oscillatory behavior at longer lag times.
In contrast, the auto-correlations from temporally shuffled surrogate data (dashed lines) remain low (within $\pm$0.2) across all lag times, confirming that the observed correlations are not artifacts of the estimation procedure.
Overall, these results demonstrate that \Model effectively captures conformational dynamics from the time series of HS-AFM images.
